\subsection{Khái niệm pipeline: stages, jobs, artifacts, cache, needs và rules}

Trong GitLab, mỗi lần lập trình viên đẩy mã nguồn lên một nhánh, hoặc khi một merge
request được tạo ra hay cập nhật, GitLab CI/CD sẽ tự động kích hoạt một \textit{pipeline}:
một luồng công việc gồm nhiều bước kiểm tra, biên dịch và triển khai chạy nối tiếp nhau.
Pipeline giúp nhóm phát hiện lỗi sớm, chuẩn hóa quy trình build và loại bỏ thao tác thủ công
dễ sai sót khi phát hành phiên bản mới cho dự án ASPX \cite{gitlabci}.

Một pipeline gồm sáu khái niệm cốt lõi mà sinh viên cần nắm vững trước khi viết tệp cấu hình:

\begin{description}
\item[Stages] là các giai đoạn lớn chạy tuần tự, ví dụ \verb|build| rồi đến \verb|test| rồi
đến \verb|deploy|. Một stage chỉ bắt đầu khi mọi job của stage trước đã thành công.
\item[Jobs] là đơn vị công việc nhỏ nhất trong pipeline, ví dụ job biên dịch solution hay
job chạy kiểm thử. Các job cùng một stage sẽ chạy song song nếu có đủ runner rảnh.
\item[Artifacts] là các tệp kết quả mà job giữ lại sau khi chạy xong, ví dụ thư mục
\verb|publish/| chứa website đã biên dịch, hoặc báo cáo kiểm thử định dạng JUnit. Các job ở
stage sau có thể tải artifacts của stage trước về để sử dụng tiếp.
\item[Cache] là cơ chế lưu trữ tạm giúp tăng tốc các lần chạy sau, ví dụ cache thư mục gói
NuGet đã tải. Khác với artifacts, cache không đảm bảo toàn vẹn và có thể bị xóa bất cứ lúc nào.
\item[Needs] cho phép một job chạy ngay khi các job mà nó phụ thuộc hoàn thành, thay vì chờ
cả stage trước kết thúc. Nhờ \verb|needs|, pipeline thực thi theo đồ thị có hướng nên nhanh hơn.
\item[Rules] quyết định job có được tạo ra hay không dựa trên tên nhánh, tag hay loại sự kiện
kích hoạt. Ví dụ job triển khai production chỉ chạy khi nhánh là \verb|main|.
\end{description}

Hình \ref{fig:c52-pipeline} minh họa pipeline ba giai đoạn điển hình cho dự án ASPX: mã nguồn
sau khi commit được build bằng MSBuild, kiểm thử bằng VSTest, rồi triển khai lên staging.
\begin{figure}[H]
\centering
\begin{tikzpicture}[>=Stealth, thick, box/.style={draw, rounded corners, fill=blue!10, minimum width=2.1cm, minimum height=0.9cm, align=center, font=\small}]
\node[box] (commit) {Commit};
\node[box, right=0.8cm of commit] (build) {Build};
\node[box, right=0.8cm of build] (test) {Test};
\node[box, right=0.8cm of test] (deploy) {Deploy};
\draw[->] (commit) -- node[midway, above, font=\footnotesize] {push} (build);
\draw[->] (build) -- node[midway, above, font=\footnotesize] {msbuild} (test);
\draw[->] (test) -- node[midway, above, font=\footnotesize] {vstest} (deploy);
\end{tikzpicture}
\caption{Luồng pipeline ba giai đoạn cho dự án ASPX}
\label{fig:c52-pipeline}
\end{figure}

\subsection{Tệp cấu hình hoàn chỉnh cho dự án ASPX 4.8.1}

Mọi pipeline của GitLab đều được khai báo trong tệp \verb|.gitlab-ci.yml| đặt ở thư mục gốc
của repository. Liệt kê \ref{lst:c52-full-yml} là một tệp hoàn chỉnh, áp dụng trực tiếp cho
solution \verb|AspxDemo_NET481.sln| gồm một project Web Forms và một project kiểm thử đơn vị. Ví dụ giả
định runner Windows dùng shell PowerShell là mặc định của executor \verb|shell| trên Windows.
\begin{lstlisting}[language=yaml,caption={Tệp cấu hình GitLab CI đầy đủ cho ASPX 4.8.1},label={lst:c52-full-yml}]
# Full GitLab CI pipeline for ASPX (.NET Framework 4.8.1)
# Requires a Windows runner with MSBuild, NuGet and VSTest
stages:
  - build
  - test
  - deploy
variables:
  BUILD_CONFIG: "Release"
  SOLUTION: "AspxDemo_NET481.sln"
  PUBLISH_DIR: "publish"
# Cache NuGet packages between pipelines on the same runner
cache:
  key: "nuget-$CI_COMMIT_REF_SLUG"
  paths:
    - packages/
# Stage 1: restore packages and build the solution
build:
  stage: build
  tags:
    - windows
    - msbuild
  script:
    - nuget restore $env:SOLUTION
    - msbuild $env:SOLUTION /p:Configuration=$env:BUILD_CONFIG /p:DeployOnBuild=true /p:PublishDir=$env:PUBLISH_DIR\
  artifacts:
    paths:
      - publish/
      - AspxDemo.Tests/bin/Release/
    expire_in: 1 week
# Stage 2: run unit tests with VSTest, reuse build artifacts
test:
  stage: test
  tags:
    - windows
  needs:
    - job: build
      artifacts: true
  script:
    - vstest.console.exe AspxDemo.Tests\bin\$env:BUILD_CONFIG\AspxDemo.Tests.dll /Logger:"trx;LogFileName=results.trx"
    - dotnet tool install -g trx2junit; if ($LASTEXITCODE -ne 0) { echo "trx2junit already installed" }
    - trx2junit TestResults\results.trx --output TestResults\junit.xml
  artifacts:
    when: always
    reports:
      junit: TestResults/junit.xml
    paths:
      - TestResults/
# Stage 3: deploy to staging, only on the main branch
deploy:
  stage: deploy
  tags:
    - windows
  needs:
    - build
  rules:
    - if: $CI_COMMIT_BRANCH == "main"
      when: on_success
    - when: never
  script:
    - xcopy publish \\staging-server\wwwroot\AspxDemo /E /Y
  environment:
    name: staging
    url: http://staging.example.com
\end{lstlisting}

Ý nghĩa từng khối trong tệp cấu hình trên như sau:

\begin{enumerate}
\item Khối \verb|stages| khai báo ba giai đoạn chạy tuần tự: biên dịch, kiểm thử rồi triển khai.
Thứ tự khai báo chính là thứ tự thực thi của pipeline.
\item Khối \verb|variables| định nghĩa các biến dùng chung như cấu hình build và tên solution,
giúp đổi cấu hình một chỗ mà không phải sửa từng job. Khối \verb|cache| giữ lại thư mục gói NuGet
để lần chạy sau không phải tải lại từ Internet.
\item Job \verb|build| gắn hai tag \verb|windows| và \verb|msbuild| để GitLab chọn đúng runner có
cài MSBuild. Job chạy \verb|nuget restore| rồi \verb|msbuild| với tham số xuất bản ra thư mục
\verb|publish/|, sau đó lưu thư mục này thành artifacts trong một tuần.
\item Job \verb|test| dùng \verb|needs| để chạy ngay khi job \verb|build| xong và nhận đúng bộ
artifacts của nó. Kết quả VSTest khai báo ở mục \verb|reports:junit| để GitLab hiển thị trên giao
diện merge request, kèm \verb|when: always| để luôn lưu báo cáo kể cả khi có test thất bại.
\item Job \verb|deploy| dùng \verb|rules| để chỉ chạy khi nhánh là \verb|main|; mọi nhánh khác đều
bị bỏ qua nhờ dòng \verb|when: never|. Lệnh \verb|xcopy| chép website lên máy chủ staging, đồng thời
khối \verb|environment| đăng ký môi trường để theo dõi lịch sử triển khai.
\end{enumerate}

\subsection{Runner: shared, specific, group và vai trò của tag}

Runner là chương trình trung gian chạy trên một máy tính riêng, chịu trách nhiệm nhận job từ
GitLab, thực thi từng dòng trong mục \verb|script|, rồi gửi kết quả và artifacts trả về máy chủ.
Bản thân GitLab chỉ điều phối pipeline, còn mọi lệnh build đều chạy trên runner. GitLab phân
runner thành ba loại theo phạm vi phục vụ, so sánh trong Bảng \ref{tab:c52-runner}:
\begin{table}[H]
\centering
\caption{So sánh ba loại runner trong GitLab}
\label{tab:c52-runner}
\begin{tabular}{p{1.6cm}p{2.6cm}p{2.4cm}p{3.8cm}}
\toprule
Loại & Phạm vi phục vụ & Người quản lý & Khi nào nên dùng \\
\midrule
Shared & Mọi project trên máy chủ & Quản trị viên hệ thống & Đồ án môn học và bài tập nhỏ, không cần môi trường đặc biệt \\
Group & Mọi project trong một group & Trưởng nhóm hoặc quản trị viên & Nhóm sinh viên dùng chung một runner riêng của nhóm \\
Specific & Một project duy nhất & Chủ project & Project ASPX cần runner Windows gắn với đúng một repository \\
\bottomrule
\end{tabular}
\end{table}

Mỗi runner được gắn một danh sách tag mô tả năng lực, ví dụ \verb|windows|, \verb|msbuild| hay
\verb|aspx|. Khi job khai báo mục \verb|tags|, GitLab chỉ giao job cho runner có đầy đủ các tag
tương ứng; job không khai báo tag sẽ chạy trên bất kỳ runner rảnh nào. Liệt kê
\ref{lst:c52-register} đăng ký một runner Windows mới trên máy build.
\begin{lstlisting}[language=bash,caption={Đăng ký Windows runner cho dự án ASPX},label={lst:c52-register}]
# Run on the Windows build machine as Administrator
gitlab-runner register --url https://gitlab.com --token glpat-9f8e7d6c5b4a3210 --executor shell --shell powershell --tag-list "windows,msbuild,aspx" --description "aspx-win-builder-01"
gitlab-runner install
gitlab-runner start
gitlab-runner verify
# After verify, the runner appears under Settings > CI/CD > Runners
\end{lstlisting}

Với dự án ASPX trên .NET Framework 4.8.1, nhóm bắt buộc phải dùng runner Windows có cài MSBuild
đầy đủ, vì .NET Framework chỉ chạy trên Windows còn Web Forms không thể biên dịch bằng bộ công cụ
đa nền tảng trên Linux \cite{msmsbuild}. Cụ thể, MSBuild cần các assembly tham chiếu của .NET
Framework, trình biên dịch trang \verb|.aspx| và VSTest để chạy kiểm thử; các thành phần này không
có trên runner Linux dùng image \verb|ubuntu-latest|. Vì vậy mọi job trong Liệt kê
\ref{lst:c52-full-yml} đều gắn tag \verb|windows| để tránh bị điều nhầm sang runner Linux.

\subsection{Biến CI/CD: masked, protected và cách đọc qua biến môi trường}

Pipeline thường cần những thông tin nhạy cảm như chuỗi kết nối cơ sở dữ liệu, mật khẩu triển khai
hay khóa ký bản phát hành. GitLab cho phép khai báo các giá trị này ở mục Settings, CI/CD,
Variables thay vì ghi trực tiếp vào tệp cấu hình, nhờ đó mã nguồn công khai không bao giờ chứa
bí mật thật.

Hai thuộc tính quan trọng nhất của biến CI/CD là:

\begin{description}
\item[Masked] yêu cầu GitLab che giá trị biến trong nhật ký job, thay bằng chuỗi \verb|[MASKED]|.
Thuộc tính này phù hợp với mật khẩu và khóa bí mật, nhưng giá trị phải đủ dài và không chứa ký tự
xuống dòng thì GitLab mới che được.
\item[Protected] yêu cầu biến chỉ được truyền cho pipeline chạy trên nhánh hoặc tag được bảo vệ,
ví dụ nhánh \verb|main|. Nhờ đó, merge request từ nhánh lạ không thể đánh cắp bí mật triển khai
bằng cách in biến ra nhật ký.
\end{description}

Trong tệp \verb|.gitlab-ci.yml|, biến được đọc theo cú pháp của shell trên runner: PowerShell dùng
\verb|$env:TEN_BIEN|, cmd dùng \verb|%TEN_BIEN%|, còn Bash dùng \verb|$TEN_BIEN|. Với mã nguồn C\# của
trang ASPX, cách đơn giản nhất là đọc biến môi trường của tiến trình đang chạy, đồng thời giữ phương
án dự phòng đọc từ \verb|Web.config| khi chạy trên máy lập trình viên như trong Liệt kê \ref{lst:c52-env}.
\begin{lstlisting}[language={[Sharp]C},caption={Đọc biến môi trường từ mã nguồn trang ASPX},label={lst:c52-env}]
// Read connection string from a CI/CD variable,
// fallback to Web.config on developer machines
string conn = System.Environment.GetEnvironmentVariable("DB_CONN");
if (string.IsNullOrEmpty(conn))
{
    conn = System.Configuration.ConfigurationManager
        .ConnectionStrings["Default"].ConnectionString;
}
using (var db = new System.Data.SqlClient.SqlConnection(conn))
{
    db.Open();
    // Run queries here with the opened connection
}
\end{lstlisting}

Quy tắc an toàn mà nhóm cần tuân thủ: mọi bí mật đều nằm trong Variables với cả hai cờ masked và
protected, tệp \verb|Web.config| trong repository chỉ chứa chuỗi kết nối giả cho môi trường phát triển,
và không ai được in toàn bộ biến môi trường ra nhật ký job.

\subsection{Môi trường triển khai và Review Apps}

Khối \verb|environment| trong mỗi job triển khai đăng ký một môi trường có tên và địa chỉ truy cập,
ví dụ \verb|staging| cho kiểm thử nội bộ và \verb|production| cho người dùng thật. GitLab gom lịch sử
triển khai của từng môi trường vào trang Environments, nơi nhóm xem được commit nào đang chạy ở đâu.

Liệt kê \ref{lst:c52-env-deploy} minh họa ba kiểu triển khai phổ biến: staging tự động mỗi khi nhánh
\verb|main| có commit mới, production triển khai thủ công bằng nút bấm để có người kiểm soát cuối cùng,
và Review App tự động dựng một môi trường tạm cho từng merge request để người review trải nghiệm trực
tiếp thay đổi trên trình duyệt.
\begin{lstlisting}[language=yaml,caption={Môi trường staging, production và Review App},label={lst:c52-env-deploy}]
# Staging deploys automatically on every main commit
deploy_staging:
  stage: deploy
  tags:
    - windows
  rules:
    - if: $CI_COMMIT_BRANCH == "main"
  script:
    - xcopy publish \\staging-server\wwwroot\AspxDemo /E /Y
  environment:
    name: staging
    url: http://staging.example.com
# Production needs one manual approval before deploy
deploy_production:
  stage: deploy
  tags:
    - windows
  rules:
    - if: $CI_COMMIT_BRANCH == "main"
      when: manual
  script:
    - xcopy publish \\prod-server\wwwroot\AspxDemo /E /Y
  environment:
    name: production
    url: http://shop.example.com
# Temporary Review App per merge request, removed after one day
review_app:
  stage: deploy
  rules:
    - if: $CI_MERGE_REQUEST_ID
  script:
    - echo Deploy review app for $CI_MERGE_REQUEST_ID
  environment:
    name: review/$CI_COMMIT_REF_SLUG
    url: http://$CI_ENVIRONMENT_SLUG.example.com
    on_stop: stop_review_app
    auto_stop_in: 1 day
# Manual job that stops the Review App above
stop_review_app:
  stage: deploy
  rules:
    - if: $CI_MERGE_REQUEST_ID
      when: manual
  script:
    - echo Stop review app
  environment:
    name: review/$CI_COMMIT_REF_SLUG
    action: stop
\end{lstlisting}

Khi phiên bản mới gặp sự cố trên production, GitLab hỗ trợ rollback theo hai cách: mở trang
Environments và bấm triển khai lại trên bản ổn định gần nhất để chạy lại pipeline cũ, hoặc revert
commit gây lỗi trên nhánh \verb|main| để pipeline tự động triển khai bản đã sửa. Cả hai cách đều giữ
lại đầy đủ lịch sử ai đã rollback và rollback về commit nào, giúp nhóm truy vết sự cố thay vì triển
khai thủ công không để lại dấu vết.

Tóm lại, pipeline GitLab mang lại cho đồ án ASPX một quy trình khép kín: mỗi commit trên nhánh
\verb|main| đều được build, kiểm thử và triển khai tự động lên đúng môi trường, mọi bí mật nằm trong
biến CI/CD được che và bảo vệ, còn mỗi merge request đều có Review App riêng để giảng viên và thành
viên trong nhóm dùng thử trước khi approve và merge.